\documentclass[11pt]{article}
\usepackage[margin=25mm]{geometry}
\usepackage[T1]{fontenc}
\usepackage{lmodern}
\usepackage{amsmath,amssymb,booktabs,array,graphicx}
\usepackage[colorlinks=true,linkcolor=blue,citecolor=blue,urlcolor=blue]{hyperref}
\usepackage[braket,qm]{qcircuit}
\newtheorem{theorem}{Theorem}
\newtheorem{lemma}{Lemma}
\newtheorem{proposition}{Proposition}
\newenvironment{proof}[1][Proof]{\par\noindent\textit{#1.}\ }{\hfill\textit{QED}\par}
\renewcommand{\labelitemi}{$\bullet$}
\setlength{\emergencystretch}{3em}
\newcommand{\OSR}{\operatorname{OSR}}
\newcommand{\CNOT}{\mathrm{CX}}
\newcommand{\Sqi}{\mathrm{SQiSW}}
\newcommand{\A}{\mathcal A}
\newcommand{\Id}{\mathbb I}
\title{Controlling two-qubit involutions:\\CNOT costs and symbolic invariants}
\author{Gregory Morse \and Peter Rakyta}
\date{7 November 2026}
\begin{document}
\maketitle
\begin{abstract}
We classify the ancilla-free CNOT cost of controlling a traceless Hermitian
involution on two qubits, allowing arbitrary one-qubit unitaries. The minimum
is four at operator-Schmidt rank four, three at rank two, and one or two at
rank one according to local support. A control-space rigidity argument
excludes three arbitrary-angle Ising interactions in the rank-four case.
Controlled fermionic SWAP attains this bound with a literal four-CNOT,
six-rotation circuit, whereas three unrestricted two-qubit gates suffice.
We extend the normal-form analysis to every spectral stratum of two-qubit
involutions. Trace-\(\pm2\) reflections form a one-Schmidt-angle family
joining controlled-controlled-Z and Fredkin: we give exact eight-CNOT and
six-unrestricted-gate upper bounds and partial lower bounds in its interior.
Finally, we reformulate the known two-SQiSW region as a necessary-and-sufficient
real polynomial test in matrix entries, suitable for symbolic families.
The accompanying symbolic library, phase-resolved circuits, and row-level
certificates provide reproducible constructions with explicit optimality
scopes. Detailed compiler experiments and numerical gap audits are supplied
separately.
\end{abstract}

\section{Introduction and related work}
Controlling an operator preserves more structure than its ordinary two-qubit
local-equivalence class. In particular, local factors in a target block become
conditional operations. We exploit this observation to classify a continuous
class of three-qubit gates: $C(V)=\operatorname{diag}(I_4,V)$, where
$V=V^\dagger$, $V^2=I_4$, and $\operatorname{tr}V=0$. A two-angle normal
form gives matching constructions on every operator-Schmidt-rank stratum.
A separate control-space argument excludes three Ising interactions on the
rank-four stratum. The resulting classification is the main contribution.

Hermitian unitaries are observables with eigenvalues $\pm1$; controlling one
is the central operation in a Hadamard test of its expectation value.
With a $|+\rangle$ control, the probability of the final $+1$ outcome is
$(1+\langle V\rangle)/2$. Controlled Pauli products are elementary examples;
locally conjugated observables and two-qubit reflections give less trivial
instances. Controlled oracles and reflections also occur in linear-combination
and qubitization constructions \cite{qubitization}. Our question is the exact
interaction cost of such a small controlled block. It does not require or
imply a complete algorithmic speedup. Spectral multiplicity separates the
balanced two-angle family from the one-angle rank-one-reflection family
containing ordinary Fredkin and controlled-controlled-Z.

Gate-count lower bounds based on operator-Schmidt strength originate in
work including Nielsen et al.\ \cite{nielsen}. Yu, Duan and Ying proved that
Toffoli requires five unrestricted two-qubit gates \cite{yuduan}; Yu and Ying
classified controlled-controlled gates and proved the five-gate unrestricted
Fredkin bound \cite{yuying}. Their controlled-product arguments are directly
relevant here. Shende and Markov studied block-diagonal three-qubit circuits
with fewer than six CNOTs and proved Toffoli's six-CNOT minimum
\cite{toffolicnot}. Our Ising proof uses the same general principle of
restricting intermediate control frames, with an explicit Schmidt-space
rigidity lemma. Balanced involutions exclude the nonbalanced relative
blocks of ordinary Fredkin and Toffoli; our smaller counts concern a
different class.

Two-qubit invariant and Cartan methods are established \cite{makhlin,shende}.
The universal two-B construction is due to Zhang et al.\ \cite{bgate}.
Huang et al.\ give analytical SQiSW synthesis, its two-pulse Weyl region,
and its universal three-pulse bound \cite{sqisw}. Our second result is an
real polynomial reformulation of that known region, with an exact algebraic
proof and implementation. Peterson, Crooks, and Smith characterize fixed-depth
two-qubit reachability as monodromy polytopes \cite{peterson}. The
trace-invariant small-CNOT tests of Shende, Bullock, and Markov
\cite{shende} are the closest methodological analogue. The benefit here is a
test on symbolic matrix entries without solving an eigenvalue problem; for a
numerical $4\times4$ matrix, ordinary canonicalization is already inexpensive. Our symbolic charts implement these established pulse
counts with scalar tracking and explicit degeneracy branches; the counts
and the existence of analytical local frames are prior results.

The accompanying software joins OSR factorization, native analytical charts,
and exact local rewrites. Quartz \cite{quartz} and QUESO \cite{queso} already
synthesize verified, gate-set-dependent optimizations. UniversalQCompiler
provides analytical synthesis in Mathematica \cite{universalqcompiler}.
Single-qubit exact Clifford+T synthesis and normal forms have an extensive
literature \cite{matsumoto,kmm,giles,gilesma}; the 192 phase-resolved
one-qubit Clifford circuits in our atlas are a small reference and regression
suite. Neither their number nor an easily reproducible shortest-path search
is presented as a major synthesis advance. Earlier numerical OSR partition
resynthesis \cite{partition} is a separate project and supplies no numerical
evidence for the identities proved here.

\section{Resources, phase, and operator-Schmidt rank}
All gate-count theorems use arbitrary one-qubit unitaries, all-to-all
connectivity, fixed input/output wires, arbitrary input states, and no
ancillas. An Ising interaction is locally equivalent to $e^{itZZ}$;
its angle may vary between uses. CNOT is one member of this class.
An unrestricted two-qubit gate is a different, strictly larger resource.
For supplied hardware alphabets, the archive counts the declared primitive;
a theorem in the relaxed model transfers as a lower bound, and requires a
separate native construction to prove a matching optimum.

Literal synthesis requires $U=C$. Projective synthesis requires
\begin{equation} U=e^{i\gamma}C,\qquad |e^{i\gamma}|=1,\label{eq:phase}\end{equation}
with the scalar recorded at zero physical cost. A physical GP gate is never
inserted. At equal declared cost we prefer literal equality. Controlling a
projective word requires $P(\gamma)$ on the added control, since
$C(U)=P_c(\gamma)C(C)$. The release proves a concrete controlled-Sycamore
substitution both with and without this correction. This illustrates the
contract; SDKs that track global phase can already implement this rule.

For a cut $A|B$, realign $U_{ab,a'b'}$ with row $(a,a')$ and column $(b,b')$.
The rank is $r_{A|B}(U)=\OSR_{A|B}(U)$. It obeys
\begin{equation}
 r_{A|B}(UV)\le r_{A|B}(U)r_{A|B}(V),\qquad
 r_{A|B}(L_A\otimes L_B)=1.\label{eq:cut}
\end{equation}
For a controlled operator, its \emph{control OSR} means this rank on
$c|ab$. If its two target branches are linearly independent it is two,
and its control-side Schmidt space is the full diagonal space
$\mathcal D=\operatorname{span}\{|0\rangle\langle0|,|1\rangle\langle1|\}$.
For a target cut, the $V$ block is a submatrix of the realignment of $C(V)$.
Thus $\OSR(V)=4$ forces ranks $(2,4,4)$ on the three singleton cuts.
Every crossing Ising interaction has rank at most two, so Eq.~\eqref{eq:cut}
constrains the number and placement of interactions.

\section{A rank-four controlled-block lower bound}
\begin{theorem}\label{thm:four}
If $\OSR(V)=4$, then $C(V)$ requires at least four Ising-class interactions,
even with arbitrary interaction angles, arbitrary one-qubit unitaries, and
an arbitrary scalar. Controlled fermionic SWAP attains this bound with
four CNOTs.
\end{theorem}
\subsection{Control-space rigidity}
\begin{lemma}\label{lem:rigidity}
If invertible one-qubit matrices $P,Q$ map the full diagonal matrix space
\[\mathcal D=\operatorname{span}\{|0\rangle\langle0|,
|1\rangle\langle1|\}\]
onto itself by $D\mapsto PDQ$, then $P,Q$
are monomial with the corresponding common permutation. Thus locally
equivalent controlled operators with control OSR two have corresponding
conditional target blocks, up to phases, exchange of the two branches,
and common target-local input/output frames.
\end{lemma}
\begin{proof}
The only rank-one lines in $\mathcal D$ are its two coordinate projectors:
$\det\operatorname{diag}(a,b)=ab$. An invertible left/right map preserves
rank and must permute these two lines. Its column and row factors
therefore have the stated form. Control Schmidt spaces of rank-two
controlled operators are full copies of $\mathcal D$, proving the second
claim. In particular, a locally equivalent realization of
$\operatorname{diag}(\Id,V)$ has a local conditional branch. Its relative
target block has the same OSR as $V$, also after branch inversion.
\end{proof}

\subsection{Exclusion of three interactions}
\begin{proof}[Proof of the lower bound in Theorem~\ref{thm:four}]
Call the targets $a,b$. The singleton cuts have ranks $2,4,4$ for
$c,a,b$: the rank-four assertion follows from the $V$ block, since a
target cut has maximum rank four. Each crossing Ising interaction has
rank at most two. Consequently $a$ and $b$ each have degree at least two
and $c$ at least one in an interaction graph. Two interactions cannot
satisfy these demands. With three, the only edge multisets are a triangle
$ca,cb,ab$, or two $ab$ edges and one control edge. The triangle has six time orders. The doubled-target-edge case has three placements of the control edge for each of its two possible targets, also six. These are the twelve surviving edge words.

First consider one control crossing and two internal target interactions.
Absorb local frames into adjacent gates and put the single crossing in
controlled form. Its target branches are
$W_+W_-$ and $W_+QW_-$, where $Q$ is a one-wire unitary and $W_\pm$
contain the target operations after/before the crossing. By
Lemma~\ref{lem:rigidity}, choose a local branch as the first one
(inverting $Q$ if the branches are exchanged). If both internal interactions
lie on one side, the corresponding product must be local, so the relative
branch is local. If one lies on each side, the relative branch is
$W_-^\dagger QW_-$ with $W_-$ locally Ising. To make the conjugation bound explicit, write the one-wire operator as
$q_0I+q_xX+q_yY+q_zZ$. Conjugating by $e^{itZZ}$ leaves $I$ and $Z$
unchanged, sends $X\otimes I$ to a combination of $X\otimes I$ and
$Y\otimes Z$, and sends $Y\otimes I$ to a combination of $Y\otimes I$
and $X\otimes Z$. Consequently the result is
$M_0\otimes I+M_1\otimes Z$ and has OSR at most two. Outer local frames
preserve that rank. In neither case can the relative block have rank four.

Now consider the triangle. Put each control interaction in controlled
form and combine all operations on the control between them into a
one-qubit unitary $N$. Local frames at the endpoints are harmless by
Lemma~\ref{lem:rigidity}. A $2\times2$ unitary is either monomial or has
all four entries nonzero. If $N$ is monomial, the circuit remains controlled
after a possible control permutation. Each conditional target branch
contains the sole internal $ab$ interaction, with only local factors
around it. That interaction must be entangling; otherwise the singleton
rank demands already fail. Neither branch is local, a contradiction.

If $N$ has all entries nonzero, factor the target operation $K$ between
the crossings to one side. The four control blocks then span
$\{\Id,B,A,AB\}$, where $B$ is a nontrivial one-wire unitary on one
target and $A$ is the other controlled target action conjugated through
$K$. Their span dimension equals the control OSR, namely two.
If the internal interaction is outside the crossings, $K$ is local,
and $A,B$ act nontrivially on different wires. The four tensor products
are independent, contradicting rank two. If the internal interaction
lies between them, independence of $\Id,B$ forces
$A\in\operatorname{span}\{\Id,B\}$. The transformed controlled gates
then act only on $c,b$. Apart from local factors the whole circuit is
$(\Id_c\otimes K_{ab})W_{cb}$. Its $a|cb$ rank is at most
$\OSR(K)=2$, again a contradiction. These cases exclude every
three-interaction word. The four-CNOT identity in Sec.~\ref{sec:fermionic} attains the bound for $U_F$.
\end{proof}
The supplied checker verifies the exact matrix identity, cut minors,
three-edge enumeration, Ising conjugation identities, and native-class
identifications. The argument above is an analytical proof with checked
algebraic components, not a proof-assistant formalization or a conclusion
from failed numerical optimization.


\section{Complete classification for balanced involutions}
\begin{theorem}\label{thm:balanced}
Let $V=V^\dagger$, $V^2=I_4$, and $\operatorname{tr}V=0$. The minimum
CNOT and unrestricted two-qubit-gate counts of $C(V)$ are:
\begin{center}\begin{tabular}{lrr}\toprule
Relative block & CNOT & Unrestricted\\\midrule
$\OSR(V)=4$ &4&3\\
$\OSR(V)=2$ &3&3\\
$\OSR(V)=1$, one-body Pauli up to local conjugacy &1&1\\
$\OSR(V)=1$, two-body Pauli product up to local conjugacy &2&2\\
\bottomrule\end{tabular}\end{center}
The CNOT lower bounds use Theorem~\ref{thm:four} in the rank-four row
and the unrestricted two-gate exclusion below in the rank-two row.
There is no rank-three stratum. Multiplying $V$ by an arbitrary scalar phase
does not change these interaction counts when arbitrary $U(2)$ locals are
available: the scalar is supplied by a physical one-qubit phase on $c$.
\end{theorem}

\subsection{Two-angle normal form, including degeneracies}
Expand a traceless Hermitian $V$ in the real Pauli basis:
\begin{equation}
 V=\vec a\cdot\vec\sigma\otimes I+I\otimes\vec b\cdot\vec\sigma
       +\sum_{ij}T_{ij}\sigma_i\otimes\sigma_j.
\end{equation}
Multiplying Pauli matrices and setting $V^2=I$ gives exactly
\begin{equation}\label{eq:involution}
 |a|^2+|b|^2+\|T\|_F^2=1,\quad Tb=0,\quad T^Ta=0,
 \quad\operatorname{cof}(T)=ab^T.
\end{equation}
The scalar, one-body, and two-body coefficients of $V^2$ independently
verify these equations. The supplied checker expands every coefficient.

If $a,b$ are nonzero, rotate them to the $X$ and $Z$ axes respectively.
Then the $X$ row and $Z$ column of $T$ vanish. Its remaining real $2\times2$
block has a signed singular-value decomposition under the remaining rotations
in the two perpendicular planes. Signs are retained, so only proper rotations,
implemented by local unitaries, are needed. This gives
\begin{equation}\label{eq:fourterm}
 V_0=A X_a+B Z_b+C Y_aY_b+D Z_aX_b,\quad
 A^2+B^2+C^2+D^2=1,\quad AB+CD=0.
\end{equation}
For completeness, if only $a$ is nonzero, $\operatorname{cof}(T)=0$ forces
$\operatorname{rank}T\le1$ and $T^Ta=0$ makes its nonzero left singular
axis perpendicular to $a$. The two remaining terms are put into
Eq.~\eqref{eq:fourterm} by local rotations. The case with only $b$ is
symmetric. If both vanish, $T$ has rank one and unit norm, giving a Pauli
product. If $T=0$, exactly one of $a,b$ is nonzero and has unit norm.
Thus the same normal form covers every degenerate case.

The matrix $\left(\begin{smallmatrix}A&-C\\D&B\end{smallmatrix}\right)$
has rank one and Frobenius norm one, so it is an outer product of two real
unit vectors. Write these as $(\sin\alpha,\cos\alpha)^T$ and
$(\sin\beta,\cos\beta)$. With $K=\CNOT_{a,b}$ and
$Q(t)=R_y(t)ZR_y(-t)$ this proves
\begin{equation}\label{eq:normalform}
 V_0=K\,[Q(\alpha)_aQ(\beta)_b]\,K,
\end{equation}
with coefficients
$(A,B,C,D)=(\sin\alpha\sin\beta,\cos\alpha\cos\beta,
-\sin\alpha\cos\beta,\cos\alpha\sin\beta)$.
In the Pauli coefficient matrix the four coefficients occupy different rows
and columns. Hence OSR is their number of nonzero entries, which is one,
two, or four. The relation $AB+CD=0$ excludes three.

\subsection{Matching constructions and two-gate exclusions}
Controlling Eq.~\eqref{eq:normalform} uses two outer $K$ gates and one
controlled $Q$ on each target. Since every $Q$ is a one-qubit Pauli axis,
its controlled gate is locally CZ and costs one CNOT. This proves the
four-CNOT upper bound throughout the class. Local conjugating frames for
$V$ add only one-qubit gates to $C(V)$.

On rank-two boundaries there is a stronger form $V_0=K'Q_jK'^\dagger$
with $K'$ containing one CNOT. The table specifies every angle;
$S=\operatorname{diag}(1,i)$ and matrix products act right to left.
\begin{center}\small\begin{tabular}{llll}\toprule
Edge & $t$ & $K'$ & $Q_j$\\\midrule
$\alpha=0$ & $\beta$ & $H_b\CNOT_{a,b}$ & $\cos t X_b+\sin t Z_b$\\
$\alpha=\pi/2$ & $\beta$ & $H_aS_b\CNOT_{a,b}$ & $\sin t Z_a+\cos t Y_a$\\
$\beta=0$ & $\alpha$ & $S_a^\dagger\CNOT_{b,a}$ & $\cos t Z_b+\sin t Y_b$\\
$\beta=\pi/2$ & $\alpha$ & $H_a\CNOT_{a,b}$ & $\sin t Z_a+\cos t X_a$\\
\bottomrule\end{tabular}\end{center}
Other edges follow by angle periodicity and local Pauli conjugations.
The matrix identity is $C(V_0)=K'C(Q_j)K'^\dagger$:
in \emph{time order} apply $K'^\dagger$, then $C(Q_j)$, then $K'$.
This costs three CNOTs.
The checker proves all four identities with $t$ free, including the corners.

To exclude two unrestricted interactions for nonlocal $V$, their edges must
cover all three wires. If both edges meet $c$, the controlled-product
argument of Yu and Ying \cite[Proposition 2]{yuying} says that the
conditional branches are target-local; their relative product cannot be
nonlocal. If the edges are $ca,ab$, there is a single crossing of the
control cut. Its OSR must be two. A two-qubit unitary of OSR two is locally
controlled: in its Cartan expansion, rank two leaves only one nonzero
interaction coordinate. Put this gate into controlled form. With the
$ab$ interaction entirely on one side of the crossing, the local-branch
argument of Lemma~\ref{lem:rigidity} makes the relative target branch local.
The reversed order and exchanged targets give the same conclusion.
Two gates therefore cannot implement any nonlocal relative $V$.

For any balanced $V$, its $+1$ and $-1$ eigenspaces both have dimension two.
An unrestricted two-qubit change of basis $J$ satisfies
$V=JZ_aJ^\dagger$. The circuit $J^\dagger,CZ_{ca},J$, in time order,
uses three unrestricted interactions. Together with the preceding exclusion
and Theorem~\ref{thm:four}, this proves both nonlocal rows of
Theorem~\ref{thm:balanced}. In the rank-one cases, local Hermitian unitary
factors reduce to identities or Pauli axes. One-body control costs one CX;
a two-body Pauli product costs two. One gate cannot cover its three
nontrivial singleton cuts, proving the latter minimum.

The theorem depends on more than Weyl class. The relative blocks $I$,
$Z_a$, and $Z_aZ_b$ all have the local Weyl class, whereas their controlled
costs are zero, one, and two. A classification of arbitrary controlled-$V$
by the Weyl class of $V$ alone is therefore impossible. The present theorem
classifies balanced involutions, not all controlled unitaries.

\subsection{Controlled fermionic SWAP}\label{sec:fermionic}
Let $F|00\rangle=|00\rangle$, $F|01\rangle=|10\rangle$,
$F|10\rangle=|01\rangle$, and $F|11\rangle=-|11\rangle$.
Then $F=(Z_a+Z_b+X_aX_b+Y_aY_b)/2$, so its four independent Pauli
Schmidt coefficients give $\OSR(F)=4$. Equivalently,
$F=(S^\dagger\otimes S^\dagger)\mathrm{iSWAP}$.
It is Hermitian, squares to identity, and has trace zero. Ordinary SWAP
has trace two and belongs to a different class.

Figure~\ref{fig:fermionic} is a literal identity with four CNOTs and six
$R_y$ rotations. In time order its rotations and interactions are
\begin{align*}
 &R_y(-\pi/2)_a,\ \CNOT_{a,b},\ R_y(-\pi/4)_a,\ R_y(-3\pi/4)_b,\\
 &\CNOT_{c,a},\ \CNOT_{c,b},\ R_y(\pi/4)_a,\ R_y(3\pi/4)_b,
 \ \CNOT_{a,b},\ R_y(\pi/2)_a.
\end{align*}
Exact multiplication gives $C(F)$ with no leftover scalar. Six is a
local-count upper bound, not a proved local minimum. An independently
constructed two-qubit $J$, rotating the odd-parity subspace by a real
$\pi/4$ rotation, also gives $F=J^\dagger Z_aJ$; its exact three-interaction
controlled circuit is checked separately.
\begin{figure}[ht]
\centering\resizebox{\linewidth}{!}{\scalebox{1.0}{
\Qcircuit @C=0.55em @R=0.5em @!R { \\
                \nghost{{q}_{0} :  } & \lstick{{q}_{0} :  } & \qw & \qw & \qw & \ctrl{1} & \qw & \ctrl{2} & \gate{R_z\left(\pi\right)} & \qw & \qw & \qw & \qw & \qw\\
                \nghost{{q}_{1} :  } & \lstick{{q}_{1} :  } & \gate{R_y\left(\frac{\pi}{2}\right)} & \ctrl{1} & \gate{R_y\left(\frac{\pi}{4}\right)} & \targ & \gate{R_y\left(-\,\frac{\pi}{4}\right)} & \qw & \qw & \ctrl{1} & \gate{R_y\left(\frac{\pi}{2}\right)} & \gate{R_z\left(\pi\right)} & \qw & \qw\\
                \nghost{{q}_{2} :  } & \lstick{{q}_{2} :  } & \gate{R_z\left(\pi\right)} & \targ & \gate{R_y\left(-\,\frac{3\,\pi}{4}\right)} & \qw & \qw & \targ & \gate{R_y\left(-\,\frac{\pi}{4}\right)} & \targ & \gate{R_x\left(-\,\pi\right)} & \qw & \qw & \qw\\
\\ }}
}
\caption{Literal controlled fermionic SWAP. Wire zero is $c$, and wires
one and two are $a,b$. The interaction count is optimal in the Ising model.}
\label{fig:fermionic}
\end{figure}

\section{The remaining spectral strata}\label{sec:reflections}
For a two-qubit Hermitian involution the trace is $0,\pm2$, or $\pm4$.
The balanced stratum was classified above. Trace $\pm4$ gives $V=\pm I$,
so $C(V)$ is identity or $Z_c$ and needs no interactions. The remaining
strata are reflections about a single two-qubit state, up to sign.

\begin{proposition}\label{prop:reflection}
Every trace-$\pm2$ two-qubit involution is locally conjugate to $\pm R_\theta$,
where
\begin{equation}
 R_\theta=I-2|\psi_\theta\rangle\langle\psi_\theta|,\qquad
 |\psi_\theta\rangle=\cos\theta|00\rangle+\sin\theta|11\rangle,
 \quad0\le\theta\le\pi/4.
\end{equation}
Every $C(\pm R_\theta)$ has an ancilla-free literal construction with at most
eight CNOTs, or six unrestricted two-qubit gates. For $0<\theta<\pi/4$,
\begin{equation}
 4\le E_{\CNOT}(C(R_\theta))\le8,\qquad
 3\le E_{\mathrm{arb}}(C(R_\theta))\le6.
\end{equation}
At $\theta=0$ the respective exact minima are six and five. At
$\theta=\pi/4$ the unrestricted minimum is five, and the CNOT bounds are
$5\le E_{\CNOT}\le7$. The interior bounds do not assert a cost classification.
\end{proposition}
\begin{proof}
The spectral theorem writes a trace-two involution as $I-2|\psi\rangle
\langle\psi|$. Schmidt decomposition puts $|\psi\rangle$ into the stated
form using local conjugations, with phases absorbed in the local frames.
Trace minus two changes only the sign. Controlling that sign adds $Z_c$.

Let $J_\theta=\CNOT_{a,b}R_y(2\theta)_a$. Since
$J_\theta|00\rangle=|\psi_\theta\rangle$,
\begin{equation}\label{eq:reflectionword}
 C(R_\theta)=J_\theta(X_aX_b)\,CCZ_{c,a,b}\,(X_aX_b)J_\theta^\dagger.
\end{equation}
A six-CNOT $CCZ$ word plus the two outer target CNOTs gives eight;
all identities are literal. For the unrestricted count use the five-gate
$CCZ$ word, in time order,
\[
 CS_{a,b},\quad\CNOT_{c,a},\quad CS_{a,b}^\dagger,\quad
 \CNOT_{c,a},\quad CS_{c,b}.
\]
Absorb the preceding $(X_aX_b)J_\theta^\dagger$ into its first $ab$ gate.
The final $J_\theta(X_aX_b)$ is one further unrestricted $ab$ gate,
so six suffice. The release checks both words for a free real $\theta$.

Direct realignment of $R_\theta$ gives
\begin{equation}\label{eq:reflectionminor}
 \det\mathcal R(R_\theta)=-\sin^2(2\theta)[2-\sin^2(2\theta)],
\end{equation}
which is nonzero on $(0,\pi/4]$. Theorem~\ref{thm:four} therefore excludes
three CNOTs, and the preceding general two-gate exclusion gives the
unrestricted lower bound three.
At zero, $R_0$ is a negative-control CZ; its controlled gate is locally
CCZ. Shende--Markov's six-CNOT and Yu--Duan--Ying's five-unrestricted-gate
bounds apply \cite{toffolicnot,yuduan}. At $\pi/4$, a local Pauli frame maps
$|\psi_\theta\rangle$ to the singlet. Reflection about the singlet is SWAP,
so this endpoint is locally Fredkin. Yu--Ying gives its unrestricted
minimum five \cite{yuying}; the known seven-CNOT construction supplies the
stated upper bound \cite{fredkinshallow}. No seven-CNOT lower bound is used.
\end{proof}
Thus all spectral strata have explicit normal forms and constructions.
Only the balanced stratum has a complete interaction-cost classification;
the interpolation between the Toffoli and Fredkin endpoints is a concrete
remaining optimization problem.

\section{An invariant characterization of the two-SQiSW region}
Let $\Gamma=U(YY)U^T(YY)$ and $\Delta=\det U$. Define the scalar-independent
Makhlin invariants \cite{makhlin}
\begin{equation}
 g_1=\frac{(\operatorname{tr}\Gamma)^2}{16\Delta},\qquad
 g_2=\frac{(\operatorname{tr}\Gamma)^2-\operatorname{tr}(\Gamma^2)}{4\Delta}.
\end{equation}
Use $\A(x,y,z)=\exp[i(xXX+yYY+zZZ)]$, with chamber
$\pi/4\ge x\ge y\ge|z|$, and $\Sqi=\A(\pi/8,\pi/8,0)$.
\begin{proposition}\label{thm:sqisw}
A two-qubit unitary has a projective realization with at most two SQiSW
pulses and arbitrary local unitaries if and only if
\begin{equation}\label{eq:iff}
 W=g_2+1\ge0,\qquad W^2\ge8(|g_1|+\Re g_1).
\end{equation}
This is an invariant rewriting of Huang et al.'s region $x\ge y+|z|$
\cite{sqisw}, including its boundary.
\end{proposition}
\begin{proof}
In the chamber, $g_1=(u+iv)^2$ up to the immaterial sign of $v$, where
$u=\cos2x\cos2y\cos2z\ge0$, and
$g_2=\cos4x+\cos4y+\cos4z$. Thus
$|g_1|+\Re g_1=2u^2$, and Eq.~\eqref{eq:iff} is equivalent to $W\ge4u$.
Put $X=\cos2x$, $Y=\cos2y$, $Z=\cos2|z|$. Direct factorization gives
\begin{align}
 W-4u&=2(X^2+Y^2+Z^2-2XYZ-1)\nonumber\\
 &=2[X-\cos2(y+|z|)][X-\cos2(y-|z|)].\label{eq:factor}
\end{align}
When $|z|>0$, the second factor is negative: $x\ge y>y-|z|$ and
cosine decreases on the relevant interval. The first factor is nonpositive
exactly when $x\ge y+|z|$. When $z=0$, the product becomes
$(X-Y)^2\ge0$ and the chamber already satisfies $x\ge y$.
Huang et al.'s necessary-and-sufficient chamber characterization now proves
both directions. Exact polynomial checks establish Eq.~\eqref{eq:factor}
and the identity for $(u+iv)^2$, independently of random sampling.
\end{proof}
For a directly polynomial version put $a=\Re g_1$, $b=\Im g_1$ and
$W=g_2+1$. Equation~\eqref{eq:iff} is equivalent to
\begin{equation}\label{eq:polynomial}
 W\ge0,\qquad W^2-8a\ge0,\qquad W^4-16aW^2-64b^2\ge0.
\end{equation}
Indeed $W^2-8a\ge8\sqrt{a^2+b^2}$ is equivalent to its nonnegative
left-hand side and the squared inequality. For unitary $U$,
$1/\det U=\overline{\det U}$, so these are real polynomial inequalities in
the real and imaginary matrix entries, with constant denominators.
They can be substituted into symbolic families while retaining parameter
conditions; a numerical canonicalization advantage is not claimed.

The test decides zero/one/two-pulse-region membership, not the exact minimum
within that region. A separate zero/one-pulse class test is needed to assert
that two pulses are necessary. No new universal three-pulse bound is claimed.


\section{Software and data}\label{sec:data}
The implementation combines exact OSR factorization, native analytical charts,
and local rewrites, accepting identities or compositional exact certificates.
Its $327\times32$ reference grid contains 217 named one-qubit entries
(including 192 phase-resolved Cliffords), 90 other named entries, and 20
fixed register fixtures. These are catalog entries rather than independent
nonlocal families. The supplement reports coverage by physical qubit count,
obstruction type, and matched proof scope, with a row-level optimality ledger.
The finite Clifford+T shortest-word certificates are a regression/reference
resource alongside established synthesis theory \cite{matsumoto,kmm,giles,gilesma}.

The exact controlled-phase SYC chart now uses two native pulses, and DCX
has an explicit one-TK2 word. The atlas also flags remaining shorter
\emph{numerical} witnesses in the SYC and $\sqrt{\mathrm{SWAP}}$ columns,
separately from its displayed exact circuits. Parameter-family samples are
labelled as samples. The detailed compiler comparisons, chart/OSR ablations,
parameter-reuse timings, and negative peephole experiment are in the supplement.
They measure the supplied implementation and do not establish compiler
dominance or a general runtime advantage.

Sources, exact words, proof-component checks, audit circuits, the ledger,
and both complete atlas PDFs with editable LaTeX are included in the review
bundle. Content hashes and Git revisions identify every recorded source.
The repository is \url{https://github.com/rakytap/sequential-quantum-gate-decomposer}.

\section{Scope and open questions}
The balanced classification allows arbitrary one-qubit unitaries, all-to-all
connectivity, and no ancillas. A restricted local alphabet needs its own
native construction. The Ising lower bound permits arbitrary interaction
angles but does not apply to unrestricted two-qubit gates. The trace-$\pm2$
interior has partial bounds; determining its exact costs is open.
Continuous local-cost minima and all-grid interaction optimality are also
open, and twenty shared-expression scores remain unevaluated. Numerical
witnesses supply optimization leads, not exact identities or impossibility
proofs. We infer no calibrated pulse latency or fidelity gain from local-box
counts. The analytical proofs have exact checked components rather than
proof-assistant formalizations.

\paragraph{Use of AI.}
An OpenAI Codex assistant supported implementation, calculations, experiments,
source research, and drafting. The authors are responsible for the mathematics,
software, data, and final text.
\begin{thebibliography}{99}
\raggedright
\bibitem{makhlin} Y. Makhlin, Nonlocal properties of two-qubit gates and
mixed states and optimization of quantum computations, Quantum Information
Processing 1, 243--252 (2002), \href{https://arxiv.org/abs/quant-ph/0002045}{arXiv:quant-ph/0002045}.
\bibitem{shende} V. V. Shende, S. S. Bullock, and I. L. Markov, Recognizing
small-circuit structure in two-qubit operators and timing Hamiltonians to
compute controlled-NOT gates, Physical Review A 70, 012310 (2004),
\href{https://doi.org/10.1103/PhysRevA.70.012310}{doi:10.1103/PhysRevA.70.012310}.
\bibitem{bgate} J. Zhang, J. Vala, S. Sastry, and K. B. Whaley, Minimum
construction of two-qubit quantum operations, Physical Review Letters 93,
020502 (2004), \href{https://arxiv.org/abs/quant-ph/0312193}{arXiv:quant-ph/0312193}.
\bibitem{sqisw} C. Huang et al., Quantum instruction set design for
performance, Physical Review Letters 130, 070601 (2023),
\href{https://arxiv.org/abs/2105.06074}{arXiv:2105.06074}.
\bibitem{quartz} M. Xu, Z. Li, O. Padon, S. Lin, J. Pointing, A. Hirth,
H. Ma, J. Palsberg, A. Aiken, U. A. Acar, and Z. Jia, Quartz:
Superoptimization of quantum circuits, Proceedings of PLDI (2022),
\href{https://doi.org/10.1145/3519939.3523433}{doi:10.1145/3519939.3523433}.
\bibitem{queso} A. Xu, A. Molavi, L. Pick, S. Tannu, and A. Albarghouthi,
Synthesizing quantum-circuit optimizers, Proceedings of the ACM on
Programming Languages 7, PLDI, Article 140 (2023),
\href{https://doi.org/10.1145/3591254}{doi:10.1145/3591254}.
\bibitem{universalqcompiler} R. Iten, O. Reardon-Smith, E. Malvetti,
L. Mondada, G. Pauvert, E. Redmond, R. S. Kohli, and R. Colbeck,
Introduction to UniversalQCompiler,
\href{https://arxiv.org/abs/1904.01072}{arXiv:1904.01072} (2019; revised 2021).
\bibitem{qiskitparameter} IBM Quantum, Qiskit Parameter documentation,
\url{https://quantum.cloud.ibm.com/docs/en/api/qiskit/qiskit.circuit.Parameter},
accessed 7 November 2026; executed versions are in the experiment manifest.
\bibitem{virtualz} D. C. McKay, C. J. Wood, S. Sheldon, J. M. Chow, and
J. M. Gambetta, Efficient Z gates for quantum computing, Physical Review A
96, 022330 (2017), \href{https://arxiv.org/abs/1612.00858}{arXiv:1612.00858}.
\bibitem{partition} G. Morse, M. Al Farizi, J. N\'adori, L. Hajas,
T. Kozsik, and P. Rakyta, Entanglement-aware quantum circuit optimization
through globally optimal partitioning,
\href{https://arxiv.org/abs/2611.03845}{arXiv:2611.03845} (2026).

\bibitem{nielsen} M. A. Nielsen, C. M. Dawson, J. L. Dodd, A. Gilchrist,
D. Mortimer, T. J. Osborne, M. J. Bremner, A. W. Harrow, and A. Hines,
Quantum dynamics as a physical resource, Physical Review A 67, 052301 (2003),
\href{https://arxiv.org/abs/quant-ph/0208077}{arXiv:quant-ph/0208077}.
\bibitem{yuduan} N. Yu, R. Duan, and M. Ying, Five two-qubit gates are
necessary for implementing Toffoli gate, Physical Review A 88, 010304(R)
(2013), \href{https://arxiv.org/abs/1301.3372}{arXiv:1301.3372}.
\bibitem{yuying} N. Yu and M. Ying, Optimal simulation of Deutsch gates and
the Fredkin gate, Physical Review A 91, 032302 (2015),
\href{https://doi.org/10.1103/PhysRevA.91.032302}{doi:10.1103/PhysRevA.91.032302}.
The controlled-product statement is Proposition 2.
\bibitem{toffolicnot} V. V. Shende and I. L. Markov, On the CNOT-cost of
TOFFOLI gates, Quantum Information and Computation 9, 461--486 (2009),
\href{https://arxiv.org/abs/0803.2316}{arXiv:0803.2316}.
\bibitem{matsumoto} K. Matsumoto and K. Amano, Representation of quantum
circuits with Clifford and $\pi/8$ gates (2008),
\href{https://arxiv.org/abs/0806.3834}{arXiv:0806.3834}.
\bibitem{kmm} V. Kliuchnikov, D. Maslov, and M. Mosca, Fast and efficient
exact synthesis of single-qubit unitaries generated by Clifford and T gates,
Quantum Information and Computation 13, 607--630 (2013),
\href{https://arxiv.org/abs/1206.5236}{arXiv:1206.5236}.
\bibitem{giles} B. Giles and P. Selinger, Exact synthesis of multiqubit
Clifford+T circuits, Physical Review A 87, 032332 (2013),
\href{https://arxiv.org/abs/1212.0506}{arXiv:1212.0506}.
\bibitem{gilesma} B. Giles and P. Selinger, Remarks on Matsumoto and Amano's
normal form for single-qubit Clifford+T operators (2013),
\href{https://arxiv.org/abs/1312.6584}{arXiv:1312.6584}.
\bibitem{maslov} D. Maslov, Advantages of using relative-phase Toffoli gates
with an application to multiple control Toffoli optimization, Physical Review
A 93, 022311 (2016), \href{https://arxiv.org/abs/1508.03273}{arXiv:1508.03273}.
\bibitem{qiskitsynthesis} IBM Quantum, Qiskit synthesis API,
\url{https://quantum.cloud.ibm.com/docs/en/api/qiskit/synthesis},
accessed 7 November 2026.
\bibitem{fredkinshallow} P. Cruz and B. Murta, Shallow unitary decompositions
of quantum Fredkin and Toffoli gates for connectivity-aware equivalent circuit
averaging, APL Quantum 1, 016105 (2024),
\href{https://arxiv.org/abs/2305.18128}{arXiv:2305.18128}.
\bibitem{peterson} E. C. Peterson, G. E. Crooks, and R. S. Smith,
Fixed-depth two-qubit circuits and the monodromy polytope, Quantum 4, 247
(2020), \href{https://doi.org/10.22331/q-2020-03-26-247}{doi:10.22331/q-2020-03-26-247};
\href{https://arxiv.org/abs/1904.10541v4}{corrected arXiv version v4} (2021).
\bibitem{qubitization} G. H. Low and I. L. Chuang, Hamiltonian simulation by
qubitization, Quantum 3, 163 (2019),
\href{https://doi.org/10.22331/q-2019-07-12-163}{doi:10.22331/q-2019-07-12-163}.
\bibitem{cirqsyc} The Cirq Developers, Analytical decompositions into
Sycamore gates, Cirq source module \texttt{two\_qubit\_to\_sycamore.py},
\url{https://github.com/quantumlib/Cirq/blob/main/cirq-google/cirq_google/transformers/analytical_decompositions/two_qubit_to_sycamore.py},
accessed 7 November 2026; the audit uses Cirq 1.3.0.
\bibitem{hp24} K. Huang and J. Palsberg, Compiling conditional quantum gates
without using helper qubits, Proceedings of the ACM on Programming Languages
8, PLDI, Article 206 (2024),
\href{https://doi.org/10.1145/3656436}{doi:10.1145/3656436}.
\bibitem{vale24} R. Vale et al., Decomposition of multi-controlled special
unitary single-qubit gates,
\href{https://arxiv.org/abs/2302.06377}{arXiv:2302.06377} (2023).
\bibitem{tk2} Quantinuum, pytket transform API: KAKDecomposition,
\url{https://docs.quantinuum.com/tket/api-docs/transform.html}, accessed
7 November 2026. One unrestricted-angle TK2 represents any two-qubit Cartan class.
\end{thebibliography}

\end{document}
